\documentclass[11pt]{article}
\input{packages.tex}
\DeclareMathOperator{\sign}{sign}
\DeclareMathOperator{\Imm}{Im}

% Настройка колонтитулов
\pagestyle{fancy}
\fancyhf{} % Очистить все поля
\fancyhead[L]{Чил аут ребятки, копим свагу}
\fancyhead[R]{\href{https://t.me/ceagest}{Автор}}
\fancyfoot[C]{\thepage}               % Справа внизу

\setlength{\headheight}{13.6pt}

\hypersetup{
    colorlinks=true,
    linkcolor=blue,
    citecolor=blue,
    filecolor=blue,
    urlcolor=blue,
}

\begin{document}

\begin{titlepage}
    \centering
    \vspace*{\fill}  % заполняет пространство сверху

    {\Huge Лекция 1. Связность и линейная связность.
    \par}

    \vspace*{\fill}  % заполняет пространство снизу
\end{titlepage}

\setcounter{page}{1}

\section{Связность. Линейная связность}

\begin{definition}
    Метрическое пространство $(X, d)$ называется \textbf{связным}, если его нельзя представить
    в виде $X = X_1 \sqcup X_2$, где $X_1, X_2$ --- непустые открытые множества. Здесь $"\sqcup"$ обозначает дизъюнктное объединение множеств (то есть 
    объединение непересекающихся множеств). В противном случае пространство называется \textbf{несвязным}.
\end{definition}

\begin{example}
    Приведём пример несвязного метрического пространства $(X, d)$. Положим $X = B_1(0, 0) \sqcup B_1(3, 0)$, $d$ --- 
    метрика, индуцированная из пространства $(\mathbb{R}^2, \left\lVert \; \cdot \; \right\rVert_2)$ (не волнуемся, если 
    не понятна эта формулировка, ниже будет разъяснение)
    \begin{figure}[h]
    \centering
    \includegraphics[width=0.5\linewidth]{pictures/Пример несвязного пространства.jpg}
    \end{figure}
\end{example}

\begin{lemma}
    Метрическое пространство $(X, d)$ несвязно тогда и только тогда, когда в нём существует
    хотя бы одно непустое множество, которое открыто и замкнуто одновременно.
\end{lemma}

\begin{proof}
    $(\Longrightarrow):$ Пусть метрическое пространство $(X, d)$ несвязно. Тогда имеет место 
    представление $X = X_1 \sqcup X_2$, где $X_1, X_2$ --- непустые открытые множества. Так как 
    $X_1 = X \setminus X_2$ и $X_2 = X \setminus X_1$, то из закона двойственности вытекает замкнутость
    множеств $X_1$ и $X_2$, а значит требуемое доказано.
    \\
    $(\Longleftarrow):$ Пусть в метрическом пространстве $(X, d)$ есть непустое множество $X_1$, которое 
    открыто и замкнуто одновременно. Положим $X_2 := X \setminus X_1$ и из закона двойственности 
    в силу замкнутости $X_1$ получим, что $X_2$ --- открытое множество. При этом по определению 
    $X_1 \sqcup X_2 = X$, что и требовалось. 
\end{proof}

\begin{definition}[Линейная связность]
    Метрическое пространство $(X, d)$ называется \textbf{линейно связным}, если
    \[\forall x, y \in X \; \exists \gamma \in C([a, b], X): \gamma(a) = x, \; \gamma(b) = y\]
\end{definition}

\begin{remark}
    Запись $\gamma \in C([a, b], X)$ означает, что отображение $\gamma$ непрерывно на 
    $[a, b]$ и принимает значения в $X$.
\end{remark}

\begin{figure}[h]
\centering
\includegraphics[width=0.5\linewidth]{pictures/Линейная связность.jpg}
\end{figure}

\begin{theorem}
    Из линейной связности следует связность. Обратное неверно.
\end{theorem}

\begin{proof}
    Пусть метрическое пространство $(X, d)$ линейно связно. Покажем, что оно связно.
    Предположим противное, то есть $X = X_1 \sqcup X_2$, где $X_1$ и $X_2$ --- непустые 
    открытые множества.
    \begin{figure}[h]
    \centering
    \includegraphics[width=0.5\linewidth]{pictures/Теорема о связи линейной связности и связности.jpg}
    \end{figure}
    \\
    Пусть $x_1 \in X_1$ и $x_2 \in X_2$. Так как $(X, d)$ линейно связно, то 
    \[\exists \gamma \in C([a, b], X): \gamma(a) = x_1, \; \gamma(b) = x_2\]
    Положим
    \[E_1 := \{t \in [a, b]: \gamma(t) \in X_1\}\]
    \[E_2 := \{t \in [a, b]: \gamma(t) \in X_2\}\]
    Так как всё пространство $X$ представляется объединением $X_1$ и $X_2$, то $[a, b] = E_1 \sqcup E_2$. При этом
    очевидно, что $E_1 = \gamma^{-1}(X_1)$ и $E_2 = \gamma^{-1}(X_2)$, а значит по теореме 
    \ref{th:1.2} множества $E_1$ и $E_2$ открыты в $[a, b]$. Тогда отрезок $[a, b]$ представлен 
    в виде дизъюнктного объединения двух относительно открытых в $[a, b]$ множеств, что невозможно.
    Получаем противоречие, а значит требуемое доказано.
    \\
    Теперь покажем, что обратное неверно. Действительно, приведём пример связного 
    пространства, которое не является линейно связным. Рассмотрим 
    \begin{equation*}
        f(x) = 
        \begin{cases}
            \sin \displaystyle \frac{1}{x}, & x \in (0, 1] \\
            0, & x = 0
        \end{cases}
    \end{equation*}
    Положим 
    \[E_1 = \{(x, f(x)) \in \mathbb{R}^2: x \in [0, 1]\}\]
    \[E_2 = \{(0, y) \in \mathbb{R}^2: y \in [-1, 1]\}\]
    Тогда возьмём такое пространство $(X, d)$, что $X = E_1 \cup E_2$, $d$ --- метрика, индуцированная из пространства $(\mathbb{R}^2, \left\lVert \; \cdot \; \right\rVert_2)$.
    Утверждается без доказательства (смекалистый читатель без труда сможет провести его самостоятельно), что 
    определённое таким образом метрическое пространство является связным, но не является линейно 
    связным.
\end{proof}

\begin{theorem} \label{th:1.2}
    Пусть $(X, d)$ и $(Y, \rho)$ --- метрические пространства, $f \colon X \to Y$. Следующие 
    условия эквивалентны:
    \begin{enumerate}
        \item Отображение $f$ непрерывно на $X$
        \item Для любого открытого множества $V \subset Y$ выполнено, что множество $f^{-1}(V)$ открыто в $X$.
    \end{enumerate}
\end{theorem}

\begin{proof}
    $(2 \Longrightarrow 1):$ Вспомним, что в произвольном метрическом пространстве открытый 
    шар всегда является открытым множеством, то есть
    \[\forall y \in Y \; \forall \varepsilon > 0 \hookrightarrow B_{\varepsilon}(y) \text{--- открытое множество}\]
    Рассмотрим теперь
    \[x = f^{-1}(y) \in f^{-1}(B_{\varepsilon}(y))\]
    Так как по условию $2$ множество $f^{-1}(B_{\varepsilon}(y))$ открыто, то:
    \[\exists \delta(x, \varepsilon) > 0: B_{\delta(x, \varepsilon)}(x) \subset f^{-1}(B_{\varepsilon}(y))\]
    Из этого, с учётом равенства $y = f(x)$, сразу следует, что 
    \[f(B_{\delta(x, \varepsilon)}(x)) \subset B_{\varepsilon}(f(x))\]
    Итого, с учётом всего вышесказанного, имеем:
    \[\forall \varepsilon > 0 \; \exists \delta(x, \varepsilon) > 0: \forall x' \in B_{\delta(x, \varepsilon)}(x) \hookrightarrow f(x') \in B_{\varepsilon}(f(x))\]
    Это по определению и означает, что $f$ непрерывно в точке $x \in X$. Но $y \in Y$ было выбрано 
    произвольно, а значит и непрерывность показана в произвольной точке множества $X$, что и требовалось.
    \\
    $(1 \Longrightarrow 2):$ Пусть отображение $f$ непрерывно в каждой точке множества $X$. Пусть $V \subset Y$ --- 
    произвольное открытое множество. Если $f^{-1}(V) = \varnothing$, то всё очевидно. Далее $f^{-1}(V) \neq \varnothing$. 
    Пусть выбрано произвольное $y \in V$ такое, что $f^{-1}(y) \neq \varnothing$. Так как $V$ открыто, то 
    \[\exists \varepsilon > 0: B_{\varepsilon}(y) \subset V\]
    Тогда из непрерывности $f$ в точке $x = f^{-1}(y) \in X$ получим:
    \[\exists \delta(x, \varepsilon) > 0: f(B_{\delta(x, \varepsilon)}(x)) \subset B_{\varepsilon}(y) \implies B_{\delta(x, \varepsilon)}(x) \subset f^{-1}(B_{\varepsilon}(y))\]
    В итоге в силу произвольного выбора $y \in V$ имеем:
    \[\forall x \in f^{-1}(V) \; \exists \delta(x, \varepsilon) > 0: B_{\delta(x, \varepsilon)}(x) \subset f^{-1}(B_{\varepsilon}(y)) \subset f^{-1}(V)\]
    Полученное по определению означает, что множество $f^{-1}(V)$ открыто в $X$, а значит требуемое доказано.
\end{proof}

\begin{theorem}[Продвинутая теорема Больцано-Коши о промежуточном значении]
    Пусть $(X, d)$ --- линейно связное метрическое пространство. Пусть $f \colon X \to \mathbb{R}$ непрерывно в 
    каждой точке множества $X$. Тогда 
    \[\forall A \in \Imm f \; \; \forall B \in \Imm f \; \; \forall C \in [\min \{A, B\}, \max \{A, B\}] \; \; \exists x_C \in X: f(x_C) = C\]
\end{theorem}

\begin{proof}
    Пусть $A, B \in \Imm f$ и $A < C < B$. Ясно, что:
    \begin{equation*}
        \begin{cases}
            A \in \Imm f \iff \exists x_A \in X: f(x_A) = A \\
            B \in \Imm f \iff \exists x_B \in X: f(x_B) = B 
        \end{cases}
    \end{equation*}
    При этом пространство $X$ линейно связно, а значит 
    \[\exists \gamma \in C([0, 1], X): \gamma(0) = x_A, \; \gamma(1) = x_B\]
    Рассмотрим отображение $h(t) = f(\gamma(t))$ $ \; \forall t \in [0, 1]$. Так как 
    композиция двух непрерывных отображений непрерывна, то $h \colon [0, 1] \to [A, B]$ --- непрерывное отображение.
    Более того, $h(0) = A$ и $h(1) = B$. Тогда по теореме Больцано-Коши о промежуточном
    значении для числовой функции получаем:
    \[\exists t_C \in [0, 1]: h(t_C) = f(\gamma(t_C)) = C\]
    Легко видеть, что, положив $x_C = \gamma(t_C) \in X$, получим то, что и требовалось.
\end{proof}

\begin{theorem}
    Пусть $(X_1, d_1), (X_2, d_2), (X_3, d_3)$ --- непустые метрические пространства. Пусть 
    \[f \in C(X_1, X_2), \quad g \in C(X_2, X_3), \quad h := g \circ f\]
    Тогда $h \in  C(X_1, X_3)$. 
\end{theorem}

\begin{proof}
    В силу теоремы \ref{th:1.2} достаточно показать, что для всякого открытого множества
    $V \subset X_3$ выполнено, что $h^{-1}(V)$ открыто в $X_1$. Заметим, что:
    \[h^{-1}(V) = (g \circ f)^{-1}(V) = f^{-1}(g^{-1}(V))\]
    При этом всё по той же теореме множество $g^{-1}(V)$ открыто в $X_2$, а значит и множество $f^{-1}(g^{-1}(V))$ является открытым в $X_1$.
    Увидев из вышеописанного, чему же равно $f^{-1}(g^{-1}(V))$, заключаем, что множество $h^{-1}(V)$ открыто в $X_1$, что и требовалось.
\end{proof}

\begin{definition}
    Пусть $(X, d)$ --- метрическое пространство. Пусть $S \subset X$, $S \neq \varnothing$. Тогда $S$ само может 
    рассматриваться как метрическое пространство с индуцированной метрикой $d_S$, то есть 
    \[\forall x, y \in S \hookrightarrow d_S(x, y) = d(x, y)\]
    Множество $V \subset S$ называется \textbf{относительно открытым} в $S$, если оно открыто как подмножество метрического пространства $(S, d_S)$.
\end{definition}

\begin{example}
    Положим $X = \mathbb{R}$, $S = [0, 1]$, $V = \left[0, \frac{1}{2}\right)$, $d(x, y) = \left\lvert x - y \right\rvert$. Тогда 
    множество $V$ относительно открыто в $S$, но не открыто в $X$.
\end{example}

\begin{lemma}
    Пусть $(X, d)$ --- метрическое пространство, $S \subset X$, $S$ --- метрическое пространство 
    с индуцированной метрикой из $(X, d)$. Множество $V \subset S$ является относительно открытым в $S$ тогда и 
    только тогда, когда найдётся такое открытое множество $W \subset X$, что $V = W \cap S$.
\end{lemma}

\begin{proof}
    $(\Longleftarrow):$ Пусть $W$ --- открытое в $X$ множество. Покажем, что $V = W \cap S$ относительно 
    открыто в $S$. Пусть $x \in V$. Так как $x \in W$ и $W$ открыто в $X$, то 
    \[\exists \delta > 0: B_{\delta}(x) \subset W\]
    Пересекая обе части этого включения с $S$, получим, что $B_{\delta}(x) \cap S \subset V$. Но множество 
    $B_{\delta}(x) \cap S$ является шаром в пространстве $S$ относительно индуцированной метрики, а $x \in V$ было 
    выбрано произвольно. Значит $V$ --- относительно открыто в $S$, что и требовалось.
    \\
    $(\Longrightarrow):$ Пусть множество $V \subset S$ относительно открыто в $S$. Покажем, что 
    его можно реализовать как пересечение $V = U \cap S$, где $U$ --- открытое в $X$ множество.
    Зафиксируем произвольно $x \in V$. Так как $V$ относительно открыто в $S$, то 
    \[\exists \delta_x > 0: \left(B_{\delta_x}(x) \cap S\right) \subset V\]
    Положим
    \[U := \bigcup_{x \in V} B_{\delta_x}(x)\]
    Так как в произвольном метрическом пространстве открытый шар всегда является открытым множеством, то 
    $U$ --- открыто в $X$. При этом из полученного условия выше также вытекает:
    \[U \cap S = \left(\bigcup_{x \in V} B_{\delta_x}(x)\right) \cap S = \left(\bigcup_{x \in V} B_{\delta_x}(x) \cap S \right) \subset V\]
    С другой стороны, так как объединение берётся по всем точкам $x \in V$, то тривиально выполнено 
    обратное включение:
    \[V \subset \left(\bigcup_{x \in V} B_{\delta_x}(x) \cap S \right) = U \cap S\]
    Значит $V = U \cap S$, что и требовалось.
\end{proof}

\begin{lemma}[Следствие предыдущей леммы]
    Пусть $(X, d)$ и $(Y, \rho)$ --- метрические пространства, $E \subset X$, $E \neq \varnothing$, 
    $G \subset Y$, $G \neq \varnothing$. Отображение $f \colon E \to G$ непрерывно в каждой точке 
    множества $E$ по множеству $E$ тогда и только тогда, когда $f \in C(E, G)$, если при этом 
    рассматривать $E$ и $G$ как метрические пространства с индуцированными метриками (подпространства в $X$ и $Y$ соответственно).
\end{lemma}

\begin{definition}
    Множество $E$ в метрическом пространстве $(X, d)$ связно (линейно связно), если метрическое 
    пространство $(E, d)$ связно (линейно связно).
\end{definition}

\begin{definition}
    Множество $\Omega \subset X$ в метрическом пространстве $(X, d)$ называется \textbf{областью},
    если $\Omega$ --- открытое линейно связное множество.
\end{definition}

\end{document}